%-------------------------
% Resume in LateX
% Author : Mengkun She
% License : MIT
%------------------------

\documentclass[letterpaper,11pt]{article}

\usepackage{latexsym}
\usepackage[empty]{fullpage}
\usepackage{titlesec}
\usepackage{marvosym}
\usepackage[usenames,dvipsnames]{color}
\usepackage{verbatim}
\usepackage{enumitem}
\usepackage[hidelinks]{hyperref}
\usepackage{fancyhdr}
\usepackage[english]{babel}
\usepackage{tabularx}
\input{glyphtounicode}

\pagestyle{fancy}
\fancyhf{} % Clear all header and footer fields
\fancyfoot{}
\renewcommand{\headrulewidth}{0pt}
\renewcommand{\footrulewidth}{0pt}

% Adjust margins
\addtolength{\oddsidemargin}{-0.5in}
\addtolength{\evensidemargin}{-0.5in}
\addtolength{\textwidth}{1in}
\addtolength{\topmargin}{-.5in}
\addtolength{\textheight}{1.0in}

\urlstyle{same}

\raggedbottom
\raggedright
\setlength{\tabcolsep}{0in}

% Sections formatting
\titleformat{\section}{
  \vspace{-4pt}\scshape\raggedright\large
}{}{0em}{}[\color{black}\titlerule \vspace{-5pt}]

% Ensure that generate PDF is machine readable/ATS parsable
\pdfgentounicode=1

%-------------------------
% Custom commands
\newcommand{\resumeItem}[2]{
  \item\small{
    \textbf{#1}{: #2 \vspace{-2pt}}
  }
}

% Just in case someone needs a heading that does not need to be in a list
\newcommand{\resumeHeading}[4]{
    \begin{tabular*}{0.99\textwidth}[t]{l@{\extracolsep{\fill}}r}
      \textbf{#1} & #2 \\
      \textit{\small#3} & \textit{\small #4} \\
    \end{tabular*}\vspace{-5pt}
}

\newcommand{\resumeSubheading}[4]{
  \vspace{-1pt}\item
    \begin{tabular*}{0.97\textwidth}[t]{l@{\extracolsep{\fill}}r}
      \textbf{#1} & #2 \\
      \textit{\small#3} & \textit{\small #4} \\
    \end{tabular*}\vspace{-5pt}
}

\newcommand{\resumeSubSubheading}[2]{
    \begin{tabular*}{0.97\textwidth}{l@{\extracolsep{\fill}}r}
      \textit{\small#1} & \textit{\small #2} \\
    \end{tabular*}\vspace{-5pt}
}

\newcommand{\resumeSubItem}[2]{\resumeItem{#1}{#2}\vspace{-4pt}}

\renewcommand{\labelitemii}{$\circ$}

\newcommand{\resumeSubHeadingListStart}{\begin{itemize}[leftmargin=*]}
\newcommand{\resumeSubHeadingListEnd}{\end{itemize}}
\newcommand{\resumeItemListStart}{\begin{itemize}}
\newcommand{\resumeItemListEnd}{\end{itemize}\vspace{-5pt}}

%-------------------------------------------
%%%%%%  CV STARTS HERE  %%%%%%%%%%%%%%%%%%%%%%%%%%%%


\begin{document}

%----------HEADING-----------------
\begin{tabular*}{\textwidth}{l@{\extracolsep{\fill}}r}
  \textbf{\href{https://sourabhbajaj.com/}{\Large \textbf{Mengkun She}}} & Email: \href{mailto:mengkun.she@gmail.com}{mengkun.she@gmail.com}\\
  \href{https://serenitysmk.github.io}{https://serenitysmk.github.io} & Mobil: \href{tel:+49 17630148035}{+49 17630148035} \\
\end{tabular*}


%-----------EDUCATION-----------------
\section{Bildung}
  \resumeSubHeadingListStart
  	\resumeSubheading
  	{Christian-Albrechts-Universität zu Kiel}{Kiel, Deutschland}{Promotion in Informatik}{Nov 2020 -- Mai 2025}
  	\resumeItemListStart
  		\resumeItem{Topic}{Underwater Refractive Camera Calibration and 3D Scene Reconstruction}
  		\resumeItem{Advisor \& Reference}{\href{https://www.marine.informatik.uni-kiel.de/}{Prof. Dr. -Ing. Kevin K\"oser}}
  	\resumeItemListEnd
    \resumeSubheading
      {Chongqing Universität}{Chongqing, China}
      {Master of Science in Vermessungs- und Kartierungstechnik (Geoinformatik)}{Sep 2017 -- Juni 2020}
    \resumeSubheading
      {Chongqing Universität}{Chongqing, China}
      {Bachelor of Science in Vermessungs- und Kartierungstechnik (Geoinformatik)}{Sep 2013 -- Juni 2017}
  \resumeSubHeadingListEnd


%-----------EXPERIENCE-----------------
\section{Erfahrung}
  \resumeSubHeadingListStart
  
   \resumeSubheading
  	{Christian-Albrechts-Universität zu Kiel}{Kiel, Deustchland}
  	{Computer Vision Research Assistant bei Marine Data Science Group (MDS)}{July 2023 -- Mai 2025}
  	\resumeItemListStart

		\resumeItem{Open-Source C++ Unterwasser-Mapping-Software}{Integrieren Sie die refraktiven Kameramodelle und das refraktive SfM in die beliebte C++-basierte Open-Source-3D-Rekonstruktionssoftware COLMAP. (Link zum Projekt: \href{https://cau-git.rz.uni-kiel.de/inf-ag-koeser/colmap_underwater}{gitlab:colmap-underwater})}
	  	\resumeItem{Unterwasser Neural Radiance Field}{Erlernen einer medien- und lichtunabh\"angigen Szenendarstellung durch Modellierung der an der Kamera befestigten mitbewegten Lichtquelle mit einem raum- und oberfl\"achennormalabh\"angigen neuronalen Netzwerk (Aufsatz: [1]).}
  		\resumeItem{Winzige und halbtransparente Objektrekonstruktion}{Erlernen einer volumetrischen Darstellung f\"ur winzige und halbtransparente Objekte mit Hilfe von differenzierbarem Rendering und neuronalem Radianzfeld (Aufsatz: [8]).}
  		\resumeItem{Lehrassistent - Probabilistic Robotics (Winter 2023)}{Verantwortlich f\"ur die Tutorien und die Abschlusspr\"ufung.}
  \resumeItemListEnd

    \resumeSubheading
      {GEOMAR Helmholtz-Zentrum f\"ur Ozeanforschung Kiel}{Kiel, Deustchland}
      {Computer Vision Research Assistant bei Oceanic Machine Vision Group (OMV)}{Nov 2020 -- Juni 2023}
      \resumeItemListStart
        \resumeItem{Refraktive Kameramodelle und Kalibrierung}
          {Ableitung und Implementierung neuartiger Kameramodelle f\"ur Unterwasserkamerasysteme, um geometrische Verzerrungen zu ber\"ucksichtigen, die durch Brechung induziert werden (Aufsatz [5, 7]).}
        \resumeItem{Refraktive bildbasierte 3D-Rekonstruktion}
        {Erforschung und Entwicklung eines refraktiven Structure-from- Motion (SfM) f\"ur die bildbasierte 3D-Rekonstruktion aus Unterwasserbildern (Aufsatz [2]).}
        \resumeItem{Sensorfusion mit Navigation}{Entwicklung einer lose gekoppelten Fusionsstrategie f\"ur die visuelle Navigation zur Kartierung eines großr\"aumigen Meeresbodens (Aufsatz: [3]).}
        \resumeItem{Makroobjektiv-Kamerakalibrierung}{Ableitung eines neuartigen affinen Transformationskameramodells f\"ur das Fokusstapelbild und Entwicklung eines C++-basierten Kalibrierungsansatzes f\"ur Makroobjektiv-Kamerasysteme (Aufsatz: [6]).}
      \resumeItemListEnd
      
    \resumeSubheading
      {GEOMAR Helmholtz-Zentrum f\"ur Ozeanforschung Kiel}{Kiel, Deustchland}
      {Praktikum \& Masterarbeit}{Jan 2019 -- Jan 2020}
      \resumeItemListStart
        \resumeItem{BubbleBox Projekt}
          {Entwicklung eines synchronisierten Hochgeschwindigkeits-Stereokamerasystems (80 Hz) f\"ur die industrielle Bildverarbeitung. Entwicklung einer C++-basierten Software zur Messung und Quantifizierung des Volumens freigesetzter Gasblasen mithilfe von Bildverarbeitungs- und Stereo-Vision-Techniken.}
      \resumeItemListEnd
      
      \resumeSubheading
      {Computer Vision Freiberufler}{Kiel, Deustchland}
      {Als Freiberufler arbeiten}{2021 -- 2024}
      \resumeItemListStart
      	\resumeItem{Lokalisierung des Roboterarm-Endeffektors}{Entwickeln Sie ein synchronisiertes 4-Kamera-System zur Verfolgung und Lokalisierung des Endeffektors eines Roboterarms in 3D mit AruCo Maker.}
      	\resumeItem{Intrinsische und extrinsische Kalibrierung eines Kamera-Arrays}{Entwicklung einer Kalibrierungsmethode zur Kalibrierung der intrinsischen und relativen extrinsischen Parameter eines Kamera-Arrays mit 26 Kameras f\"ur das medizinische Scannen des menschlichen K\"orpers.}
      \resumeItemListEnd
  \resumeSubHeadingListEnd


%-----------Skills-----------------
\section{F\"ahigkeiten}
  \resumeSubHeadingListStart
  	\resumeItem{Fähigkeiten}{Kamerakalibrierung; 3D-Rekonstruktion; SfM / SLAM; Sch\"atzung und Optimierung; Lernen der Repr\"asentation}
    \resumeSubItem{Programmierung \& Tools}
      {C++, CMake, Python, PyTorch, OpenCV, Ceres-Solver, Linux, Git, ROS (Grundkenntnisse), LaTex}
    \resumeSubItem{Wissenschaftliche F\"ahigkeiten}
      {Probleml\"osung; Modellierung; Wissenschaftliches Zeichnen und Schreiben.}
    \resumeSubItem{Sprachen}
      {English (flie\ss end); Deustch (Grundkenntnisse); Chinesisch (Muttersprache)}
  \resumeSubHeadingListEnd

\section{Zertifikat}
\resumeSubHeadingListStart
\resumeItem{Algorithmic Toolbox}{Ausgestellt von Coursera, Anmeldeinformations-ID: PX5W3GZW8QLW}
\resumeItem{Neural Networks and Deep Learning}{Ausgestellt von Coursera, Anmeldeinformations-ID: 6ML6ZBMNEXXU}
\resumeItem{Improving Deep Neural Networks}{Ausgestellt von Coursera, Anmeldeinformations-ID: JVRXWMCD4LHH}
\resumeSubHeadingListEnd

\section{Auszeichnungen  \& Stipendien}
\resumeSubHeadingListStart
\resumeItem{Ph.D. Stipendium}{Promotionsstipendium des China Scholarship Council (CSC, 2020 -- 2024)}
\resumeItem{Reisestipendium}{Reisestipendium f\"ur junge Wissenschaftlerinnen und Wissenschaftler der Deutsche Arbeitsgemeinschaft f\"ur Mustererkennung, DAGM 2019}
\resumeSubHeadingListEnd

  
\section{Ausgew\"ahlte Publikationen}
18 Publikationen bei Google Scholar gelistet, mit insgesamt 199 Zitaten und einem H-Index von 8. Nachfolgend finden Sie ausgew\"ahlte Publikationen:
 \resumeSubHeadingListStart
 \resumeItem{[1]}{\textbf{M. She}, F. Seegr\"aber, D. Nakath, P. Sch\"ontag, K. K\"oser. Relative Illumination Fields: Learning Medium and Light Independent Underwater Scenes. (Submitted to \textit{CVPR, 2025})}
 \resumeItem{[2]}{\textbf{M. She}, F. Seegr\"aber, D. Nakath, K. K\"oser. Refractive COLMAP: Refractive Structure-from-Motion Revisited. In \textit{IROS, 2024} (\textbf{Oral})}
 \resumeItem{[3]}{\textbf{M. She}, Y. Song, D. Nakath, K. K\"oser. Semihierarchical Reconstruction and Weak-area Revisiting for Robotic Visual Seafloor Mapping. In \textit{Journal of Field Robotics}}
 \resumeItem{[4]}{\textbf{M. She}, T. Wei\ss, Y. Song, P. Urban, J. Greinert, K. K\"oser. Marine Bubble Flow Quantification Using Wide-baseline Stereo Photogrammetry. In \textit{ISPRS Photogrammetry and Remote Sensing}}
 \resumeItem{[5]}{\textbf{M. She}, D. Nakath, Y.Song, K. K\"oser. Refractive Geometry on Underwater Domes. In \textit{ISPRS Photogrammetry and Remote Sensing}}
 \resumeItem{[6]}{X. Weng*, \textbf{M. She}*, D. Nakath, K. K\"oser (*\textit{Equal Contribution}). Macal - Macro Lens Calibration and the Focus Stack Camera Model. In \textit{3DV, 2021} (\textbf{Oral})}
 \resumeItem{[7]}{\textbf{M. She}, Y. Song, J Mohrmann, K. K\"oser. Adjustment and Calibration of Dome Port Camera Systems for Underwater Vision. In \textit{GCPR, 2019} (\textbf{Oral})}
 \resumeItem{[8]}{D. Nakath, X. Weng, \textbf{M. She}, K. K\"oser. Visual Tomography: Physically Faithful Volumetric Models of Partially Translucent Objects. In 3DV, 2024}
 \resumeSubHeadingListEnd

%
%--------SKILLS------------
%\section{Skills}
%  \resumeSubHeadingListStart
%    \item{
%      \textbf{Languages}{: Scala, Python, JavaScript, C++, SQL, Java}
%      \hfill
%      \textbf{Technologies}{: AWS, Play, React, Kafka, GCE}
%    }
%  \resumeSubHeadingListEnd


%-------------------------------------------
\end{document}
